\chapter{Complexity — Padding}
%%% generated by the blueprint pipeline from TCSlib.Complexity.PolyHierarchy.Padding
\section{Overview}

Declarations of the Lean module \texttt{TCSlib.Complexity.PolyHierarchy.Padding}.

\section{Declarations}

\begin{definition}[Length-preserving padded decoder]
\lean{Complexity.PolyHierarchy.padDecode}
\label{Complexity.PolyHierarchy.padDecode}
\leanok
\uses{Turing.splitAtLastTrue}
Given two finite Boolean lists $s$ and $v$, let $u$ be the prefix returned by splitting
$v$ at its last true entry, when such an entry exists. The decoder returns $u$ if its
length equals the length of $s$, and returns $s$ otherwise; in particular, it returns
$s$ when $v$ has no true entry.
\end{definition}

\begin{theorem}[Length preservation by padded decoding]
\lean{Complexity.PolyHierarchy.length_padDecode}
\label{Complexity.PolyHierarchy.length_padDecode}
\leanok
\uses{Complexity.PolyHierarchy.padDecode}
\difficulty{4}
For any Boolean words $s$ and $v$, the padded decoding of $v$ relative to $s$ has the
same length as $s$.
\end{theorem}

\begin{theorem}[Dropping a false prefix]
\lean{Complexity.PolyHierarchy.dropWhile_replicate_false}
\label{Complexity.PolyHierarchy.dropWhile_replicate_false}
\leanok
\difficulty{4}
For every natural number $k$ and Boolean word $r$, deleting the initial false entries
from the word consisting of $k$ false entries followed by true and then $r$ gives the
word consisting of true followed by $r$.
\end{theorem}

\begin{theorem}[Splitting at a marked final true bit]
\lean{Complexity.PolyHierarchy.splitAtLastTrue_marker}
\label{Complexity.PolyHierarchy.splitAtLastTrue_marker}
\leanok
\uses{Complexity.PolyHierarchy.dropWhile_replicate_false, Turing.splitAtLastTrue}
\difficulty{3}
For every Boolean word $u$ and every $k\in\mathbb N$, splitting the word formed by
appending a true bit and then $k$ false bits to $u$ at its last true bit returns $u$.
\end{theorem}

\begin{theorem}[Decoding a marked padded word]
\lean{Complexity.PolyHierarchy.padDecode_marker}
\label{Complexity.PolyHierarchy.padDecode_marker}
\leanok
\uses{Complexity.PolyHierarchy.padDecode, Complexity.PolyHierarchy.splitAtLastTrue_marker}
\difficulty{1}
Let $s$ and $u$ be Boolean words of equal length, and let $k$ be a natural number. The
padded decoding of the word formed by appending true and then $k$ false entries to $u$,
relative to $s$, is $u$.
\end{theorem}

\begin{theorem}[A padded word decoding to a target]
\lean{Complexity.PolyHierarchy.exists_padDecode_eq}
\label{Complexity.PolyHierarchy.exists_padDecode_eq}
\leanok
\uses{Complexity.PolyHierarchy.padDecode, Complexity.PolyHierarchy.padDecode_marker}
\difficulty{3}
Let $s$ and $u$ be Boolean words of equal length, and let $B$ be a natural number
greater than their common length. There exists a Boolean word $v$ of length $B$ whose
padded decoding relative to $s$ is $u$.
\end{theorem}

\begin{theorem}[Polynomial-time padding decoder]
\lean{Complexity.PolyHierarchy.polyTimeComputable_padDecode}
\label{Complexity.PolyHierarchy.polyTimeComputable_padDecode}
\leanok
\uses{Complexity.PolyHierarchy.padDecode, Complexity.PolyTimeComputable, Complexity.PolyTimeComputable.comp, Complexity.pairFstD, Complexity.pairSndD, Complexity.polyTimeComputable_ite, Complexity.polyTimeComputable_lenEq, Complexity.polyTimeComputable_of_linear, Complexity.polyTimeComputable_pairFstD, Complexity.polyTimeComputable_pairSndD, Turing.FinTM.computesFunInTime_pairValid, Turing.FinTM.computesFunInTime_stripLast, Turing.pairDecode, Turing.pairDecode_pairEncode, Turing.pairEncode, Turing.splitAtLastTrue}
\difficulty{5}
The function on bit strings that sends $z$ to $\operatorname{pad}(\pi_1(z),\pi_2(z))$ is
polynomial-time computable, where $\pi_1,\pi_2$ are the total projections associated
with the self-delimiting pair encoding. Here $\operatorname{pad}(s,v)$ returns the
prefix before the last true bit of $v$ when that prefix exists and has length $|s|$, and
returns $s$ otherwise.
\end{theorem}

\begin{definition}[Recursive tuple decoder]
\lean{Complexity.PolyHierarchy.tupleDecode}
\label{Complexity.PolyHierarchy.tupleDecode}
\leanok
\uses{Complexity.pairFstD, Complexity.pairSndD, Turing.pairEncode}
Given two functions $\mathrm{base}$ and $\mathrm{dec}$ from pairs of bit strings to bit
strings, define a decoder at each natural-number level. At level zero it returns
$\mathrm{base}(s,w)$. At level $n+1$, it pairs the level-$n$ decoding of $s$ and the
default first projection of $z$ with $\mathrm{dec}$ applied to $s$ and the default
second projection of $z$.
\end{definition}

\begin{theorem}[Tuple decoding of an encoded pair]
\lean{Complexity.PolyHierarchy.tupleDecode_succ_pairEncode}
\label{Complexity.PolyHierarchy.tupleDecode_succ_pairEncode}
\leanok
\uses{Complexity.PolyHierarchy.tupleDecode, Turing.pairEncode}
\difficulty{1}
Let $\mathrm{base}$ and $\mathrm{dec}$ be functions from pairs of bit strings to bit
strings. For every natural number $n$ and bit strings $s,w,v$, the level-$(n+1)$ tuple
decoding of $\langle w,v\rangle$ with parameter $s$ is $\langle
u,\mathrm{dec}(s,v)\rangle$, where $u$ is the level-$n$ tuple decoding of $w$ with
parameter $s$.
\end{theorem}

\begin{theorem}[Polynomial-Time Recursive Pair Decoding]
\lean{Complexity.PolyHierarchy.polyTimeComputable_tupleDecode}
\label{Complexity.PolyHierarchy.polyTimeComputable_tupleDecode}
\leanok
\uses{Complexity.PolyHierarchy.tupleDecode, Complexity.PolyTimeComputable, Complexity.PolyTimeComputable.comp, Complexity.PolyTimeComputable.pairEncode, Complexity.pairFstD, Complexity.pairSndD, Complexity.polyTimeComputable_pairFstD, Complexity.polyTimeComputable_pairSndD, Turing.pairEncode}
\difficulty{4}
Let $b,d:\{0,1\}^*\times\{0,1\}^*\to\{0,1\}^*$ be functions such that the maps sending a
bit string $p$ to $b(\pi_1(p),\pi_2(p))$ and to $d(\pi_1(p),\pi_2(p))$ are
polynomial-time computable, where $\pi_1,\pi_2$ are the total projections associated
with the self-delimiting pair encoding. For every $n\in\mathbb{N}$, the map obtained by
starting at level zero with $p\mapsto b(\pi_1(p),\pi_2(p))$ and, at each successive
level, pairing the preceding-level value on $\pi_1(p)$ with $d(\pi_1(p),\pi_2(p))$ is
polynomial-time computable.
\end{theorem}

\begin{definition}[Unary polynomial-time computability]
\lean{Complexity.PolyHierarchy.UnaryPT}
\label{Complexity.PolyHierarchy.UnaryPT}
\leanok
\uses{Complexity.PolyTimeComputable}
A length function $\ell$ on bit strings is unary polynomial-time computable if the
function sending each bit string $y$ to the unary string consisting of $\ell(y)$ copies
of $1$ is polynomial-time computable.
\end{definition}

\begin{theorem}[Polynomial Bound for Unary Polynomial-Time Functions]
\lean{Complexity.PolyHierarchy.UnaryPT.bound}
\label{Complexity.PolyHierarchy.UnaryPT.bound}
\leanok
\uses{Complexity.PolyHierarchy.UnaryPT, Complexity.PolyTimeComputable.output_length_le}
\difficulty{3}
If $\ell$ is a function on finite binary strings whose unary representation is
computable in polynomial time, then there are natural numbers $A,a$ such that, for every
binary string $y$, $\ell(y)\le A(|y|+1)^a$.
\end{theorem}

\begin{theorem}[Polynomial-time unary output]
\lean{Complexity.PolyHierarchy.unaryPT_poly}
\label{Complexity.PolyHierarchy.unaryPT_poly}
\leanok
\uses{Complexity.PolyHierarchy.UnaryPT, Complexity.PolyTimeComputable, Complexity.PolyTimeComputable.comp, Turing.FinTM.computesFunInTime_polyUnary}
\difficulty{4}
Let $f$ be a polynomial-time computable function on bit strings, and let
$C,c\in\mathbb{N}$. The function that maps each bit string $y$ to a string of
$C\bigl(|f(y)|+1\bigr)^c$ copies of the bit $1$ is polynomial-time computable.
\end{theorem}

\begin{theorem}[Constant unary output is polynomial-time computable]
\lean{Complexity.PolyHierarchy.unaryPT_const}
\label{Complexity.PolyHierarchy.unaryPT_const}
\leanok
\uses{Complexity.PolyHierarchy.UnaryPT, Complexity.polyTimeComputable_const}
\difficulty{1}
For every $k\in\mathbb N$, the function that sends each bit string to the word
consisting of $k$ true bits is polynomial-time computable.
\end{theorem}

\begin{theorem}[Addition of Unary Polynomial-Time Functions]
\lean{Complexity.PolyHierarchy.UnaryPT.add}
\label{Complexity.PolyHierarchy.UnaryPT.add}
\leanok
\uses{Complexity.PolyHierarchy.UnaryPT, Complexity.PolyTimeComputable.comp, Complexity.PolyTimeComputable.pairEncode, Complexity.polyTimeComputable_pairConcat}
\difficulty{3}
If $\ell_1$ and $\ell_2$ are functions on finite binary strings whose unary
representations, consisting of that many $1$ bits, are computable in polynomial time,
then the function $y\mapsto \ell_1(y)+\ell_2(y)$ has a polynomial-time computable unary
representation as well.
\end{theorem}

\begin{theorem}[Composition of unary polynomial-time functions]
\lean{Complexity.PolyHierarchy.UnaryPT.comp}
\label{Complexity.PolyHierarchy.UnaryPT.comp}
\leanok
\uses{Complexity.PolyHierarchy.UnaryPT, Complexity.PolyTimeComputable, Complexity.PolyTimeComputable.comp}
\difficulty{1}
Let $\ell$ map bit strings to natural numbers, and let $f$ map bit strings to bit
strings. If the unary word function $y\mapsto 1^{\ell(y)}$ is polynomial-time computable
and $f$ is polynomial-time computable, then $y\mapsto 1^{\ell(f(y))}$ is polynomial-time
computable.
\end{theorem}
